% Section-only file: % Section-only file: % Section-only file: % Section-only file: \input{new_metrics}
% Requires amsmath; add new_metrics to the parent document's BibTeX databases.
% Definitions describe diagnostics, not empirical results or validated controls.
\section{Temporal Routing and Reasoning-Class Dynamics}
\label{sec:new-routing-metrics}

We extend aggregate routing statistics with diagnostics of allocation persistence,
changes in mixture weights, and the organization of reasoning episodes.
The definitions below separate published foundations from adaptations introduced
for this study. They measure observed behavior; a large or small value does not,
by itself, identify useful reasoning or a beneficial intervention.

\subsection{Notation and observation scope}

At routed layer $\ell$ and reasoning-token position $t$, let
$S_{\ell t}\subseteq\{1,\ldots,E_\ell\}$ be the native selected expert set,
with $|S_{\ell t}|=k_\ell$. Let $p_{\ell t}$ denote the full router probability
vector and $a_{\ell te}\geq0$ the weight actually assigned to a selected expert.
The normalized executed mixture is
\begin{equation}
 m_{\ell te}=
 \frac{\mathbf{1}\{e\in S_{\ell t}\}\,a_{\ell te}}
      {\sum_{j\in S_{\ell t}}a_{\ell tj}}.
 \label{eq:new-executed-mixture}
\end{equation}
Thus $p_{\ell t}$ describes router preferences, whereas $m_{\ell t}$ describes
the relative weights of the selected routed experts. Neither is reconstructed
from scalar confidence summaries. Shared experts are excluded from $S_{\ell t}$
and these routed-expert distributions.

We use a trailing window $W_t^{(w)}=\{t-w+1,\ldots,t\}$ contained in one
contiguous reasoning segment. The current configuration uses
$w\in\{64,256\}$, stride $64$, and lags $h\in\{1,4,16,64\}$.
Prompt tokens, padding and gaps are excluded; insufficient history yields an
unavailable metric. All logarithms are natural, so entropies and divergences are
in nats. Layer indices are preserved throughout.

\subsection{Expert-set persistence and turnover}

Let $V_{t,h}^{(w)}$ contain positions $u\in W_t^{(w)}$ for which $u-h$ is
observed in the same reasoning segment. Lagged overlap is
\begin{equation}
 O_{\ell,h}(t;w)=\frac{1}{|V_{t,h}^{(w)}|}
 \sum_{u\in V_{t,h}^{(w)}}
 \frac{|S_{\ell u}\cap S_{\ell,u-h}|}{k_\ell},
 \qquad U_\ell(t;w)=1-O_{\ell,1}(t;w).
 \label{eq:new-overlap}
\end{equation}
Overlap lies in $[0,1]$ and scores persistence of selected identities; turnover
$U$ scores adjacent replacement. Rank changes inside an unchanged set do not
increase turnover. The temporal lag construction extends the existing adjacent
overlap metric.

\paragraph{Uniform-chance correction.}
Following the normalized overlap of Olson et al.~\cite{newmetrics_olson2025},
but applying it to temporal pairs, define
\begin{equation}
 O^{\mathrm{unif}}_{\ell,h}
 =\frac{O_{\ell,h}-k_\ell/E_\ell}{1-k_\ell/E_\ell}.
 \label{eq:new-corrected-overlap}
\end{equation}
Zero denotes the uniform independent-selection baseline; one denotes identical
sets. Negative values are possible, and the expression is undefined when
$k_\ell=E_\ell$. This affine transform is not independent evidence from raw overlap.

\paragraph{Usage-preserving order control.}
Our additional control permutes routing sets within the observed trailing
$w+h$-token pool, clipped at the segment start. If the pool contains $n$ tokens
and expert $e$ appears at $n_e$ positions, the exact expected overlap under
sampling distinct token positions is
\begin{equation}
 B_{\ell,h}=\frac{\sum_e n_e(n_e-1)}{n(n-1)k_\ell},
 \qquad O^{\mathrm{excess}}_{\ell,h}=O_{\ell,h}-B_{\ell,h}.
 \label{eq:new-shuffled-overlap}
\end{equation}
This separates temporal persistence from high overlap caused simply by repeatedly
using a small expert pool. Positive excess indicates more persistence than the
same usage frequencies predict under random ordering.

\subsection{Allocation concentration, diversity and drift}

\paragraph{Local and marginal entropy.}
LIMoE distinguishes token-local entropy from entropy of the average routing
distribution~\cite{newmetrics_mustafa2022}. We adapt this distinction from
modality-level training regularization to descriptive trailing windows.
For $r\in\{p,m\}$, set $\bar r_{\ell,W}=w^{-1}\sum_{u\in W}r_{\ell u}$ and
\begin{align}
 H^{r}_{\mathrm{local}}(W)
   &=\frac{1}{w}\sum_{u\in W}H(r_{\ell u}),
 &H^{r}_{\mathrm{marginal}}(W)&=H(\bar r_{\ell,W}),\nonumber\\
 H(v)&=-\sum_e v_e\log v_e,
 &I^{r}_{\mathrm{alloc}}(W)&=H^{r}_{\mathrm{marginal}}(W)
                           -H^{r}_{\mathrm{local}}(W).
 \label{eq:new-entropies}
\end{align}
Local entropy scores within-token dispersion. Marginal entropy scores diversity
across the window. Their nonnegative difference scores how much distributions
vary across positions: it is the mutual information between a uniformly sampled
window position and an expert drawn from that position's distribution. It is
invariant to token order, so it does not itself measure temporal drift.

For native selection counts $c_e=\sum_{u\in W}\mathbf{1}\{e\in S_{\ell u}\}$,
we additionally report $H_{\mathrm{selection}}=H((c_e/(wk_\ell))_e)$.
This describes usage diversity without using mixture weights. The corresponding
selection rates $c_e/w$ sum to $k_\ell$, not one.

\paragraph{Effective expert count.}
We express marginal executed-mixture entropy as
\begin{equation}
 N_{\mathrm{eff},\ell}(W)=\exp\!\left(H^{m}_{\mathrm{marginal}}(W)\right).
 \label{eq:new-effective-experts}
\end{equation}
It is the number of equally weighted experts with the same window-level entropy,
between $1$ and $E_\ell$. It can exceed $k_\ell$ because different tokens select
different experts; it is not the number executed per token.

\paragraph{Within-window distribution drift.}
Using Jensen--Shannon divergence~\cite{newmetrics_lin1991}, split $W$ into equal
chronological halves $W^-$ and $W^+$ and define
\begin{align}
 D^r_\ell(W)&=\operatorname{JSD}(\bar r_{\ell,W^-},\bar r_{\ell,W^+}),\nonumber\\
 \operatorname{JSD}(a,b)&=H\!\left(\frac{a+b}{2}\right)
                        -\frac{H(a)+H(b)}{2}.
 \label{eq:new-jsd}
\end{align}
The score lies in $[0,\log2]$ and measures redistribution between the two halves.
Full-router and executed-mixture drift are reported separately. Our temporal
control averages the same statistic over $B=16$ random permutations of positions
inside $W$, giving $D^{r,\mathrm{shuffle}}$ and
$D^{r,\mathrm{excess}}=D^r-D^{r,\mathrm{shuffle}}$.
Excess drift can be negative. All permutations use only the observed window.

\subsection{Native selection-boundary diagnostics}

Let $s_{\ell te}$ be the scores actually used for deterministic top-$k$
selection, including any relevant selection bias. We extend the existing
boundary-margin diagnostic with a temporal summary:
\begin{equation}
 g_{\ell t}=\min_{e\in S_{\ell t}}s_{\ell te}
            -\max_{e\notin S_{\ell t}}s_{\ell te},
 \qquad \bar g_\ell(W)=\frac{1}{w}\sum_{u\in W}g_{\ell u}.
 \label{eq:new-boundary}
\end{equation}
A small gap indicates a nearly tied selected/unselected boundary in score space;
it does not measure answer robustness. Ties give zero. This metric requires
$k_\ell<E_\ell$ and verified native score semantics; unsupported routing rules
are marked unavailable.

The study-specific weak-boundary fraction is
\begin{equation}
 F^{\mathrm{weak}}_\ell(W)=\frac{1}{w}
 \sum_{u\in W}\mathbf{1}\{g_{\ell u}\leq\theta_\ell\},
 \qquad \theta_\ell=Q_{0.10}(g_\ell;\mathcal D_{\mathrm{train}}).
 \label{eq:new-weak-boundary}
\end{equation}
The threshold is fitted only on calibration data or the training fold, counting
each token once despite overlapping windows. A larger fraction means more
near-boundary selections relative to that reference. Until such a threshold is
fitted, the fraction remains unavailable; the compact correlation report uses
the mean gap only.

\subsection{Reasoning-class dynamics}

ThinkARM motivates describing reasoning through episode sequences and transition
patterns~\cite{newmetrics_li2026}. We retain the project's seven labels
$\mathcal C=\{\mathrm{Read},\mathrm{Analyze},\mathrm{Plan},\mathrm{Implement},
\mathrm{Explore},\mathrm{Verify},\mathrm{Monitor}\}$ rather than ThinkARM's
eight-class taxonomy. Our censoring and hazard diagnostics are study-specific
extensions, not claimed reproductions of ThinkARM metrics.
Let $C_i$ be the label of reasoning sentence $i$ and $n_i$ its aligned token count.
Unknown labels break sequences rather than becoming an eighth class.

\paragraph{Occupancy, transitions and transition uncertainty.}
Sentence and token shares are respectively
$f_c=\sum_i\mathbf{1}\{C_i=c\}/N$ and
$f_c^{\mathrm{tok}}=\sum_i n_i\mathbf{1}\{C_i=c\}/\sum_i n_i$,
with denominators restricted to labeled observations and label coverage reported
separately. For adjacent labeled sentences in one contiguous segment, let $N_{cd}$
count transitions from $c$ to $d$. Then
\begin{equation}
 P_{cd}=\frac{N_{cd}}{\sum_{b\in\mathcal C}N_{cb}},\qquad
 P_{\mathrm{stay}}(c)=P_{cc},\qquad
 P_{\mathrm{exit}}(c)=1-P_{cc},\qquad
 H_{\mathrm{next}}(c)=-\sum_d P_{cd}\log P_{cd}.
 \label{eq:new-transitions}
\end{equation}
These score destination preference, persistence, switching, and uncertainty over
observed next classes. Zero source support gives an undefined rate. Natural
termination is outside this seven-class matrix; hence $1-P_{cc}$ is conditional
on an observed next labeled sentence, not an unconditional stopping probability.
Conditioning additionally on $C_{i-1}$ describes second-order transition patterns.

\paragraph{Runs, returns and motifs.}
A run is a maximal contiguous sequence of the same class. Its sentence/token
length scores sustained allocation to that reasoning activity; its current age
uses only the observed prefix. Entry after a missing-label gap is unobserved, so
full duration is unavailable. A gap, cap or unknown ending right-censors the run;
observed natural termination is a competing exit event.
Return intervals measure the sentence/token distance between successive run
starts of the same class, without crossing gaps. Counts of contiguous label
bigrams and trigrams describe recurring local motifs. Motif counts alone are
not significance scores and do not establish purposeful self-correction.

\paragraph{Destination-specific exit hazards.}
For runs with observed entry, let $Y_c(a)$ count runs whose outgoing boundary at
sentence age $a$ is observable, and $D_{cd}(a)$ count exits to destination $d$ at
that age. Our discrete hazard is
\begin{equation}
 h_{c\rightarrow d}(a)=D_{cd}(a)/Y_c(a),
 \qquad d\in(\mathcal C\setminus\{c\})\cup\{\mathrm{natural\ termination}\}.
 \label{eq:new-hazard}
\end{equation}
It scores the chance of a particular observed exit at that age, conditional on
reaching an observable boundary. Right-censored runs contribute earlier
boundaries but not their unobserved final boundary. Empty risk sets remain
undefined; observed run lengths must not be mistaken for uncensored dwell times.

\subsection{Relating routing to reasoning classes}

\paragraph{Within-question class contrasts.}
For a scalar window metric $z$, let $\bar z_{q,c}$ average windows ending in
class $c$ within question $q$, and $\bar z_{q,\neg c}$ average other labeled
endpoints. Our contrast is
\begin{equation}
 \Delta_c(z)=\frac{1}{|\mathcal Q_c|}
 \sum_{q\in\mathcal Q_c}(\bar z_{q,c}-\bar z_{q,\neg c}),
 \label{eq:new-class-contrast}
\end{equation}
where $\mathcal Q_c$ contains questions with both kinds of window. Layers are
averaged equally within each question. Positive values indicate larger metric
values around that class. A trailing window can include earlier classes, so this
is an endpoint-conditioned association. The saved original class-view comparison
instead uses class minus all reasoning and is kept as a distinct estimand.

\paragraph{Expert-specific association diagnostics.}
The broader diagnostic interface also compares expert activation rates, adapting
SteerMoE's contrast principle~\cite{newmetrics_fayyaz2026} from paired behavioral
prompts to observed classes. For tokens with label $c$, define
\begin{equation}
 R_{\ell e\mid c}=\frac{\sum_t\mathbf{1}\{C(t)=c,\ e\in S_{\ell t}\}}
                          {\sum_t\mathbf{1}\{C(t)=c\}},
 \qquad A_{\ell e,c}=R_{\ell e\mid c}-R_{\ell e\mid\neg c}.
 \label{eq:new-activation-contrast}
\end{equation}
This is a token-weighted selection-rate difference; positive values indicate
class-associated recruitment. Expert identity always includes its layer.
A secondary normalized pointwise mutual information score, following the
association measure used by RICE~\cite{newmetrics_wang2025}, is
\begin{equation}
 \operatorname{nPMI}(X_{\ell e},c)=
 \frac{\log\!\left[P(X_{\ell e}=1,C=c)/
                (P(X_{\ell e}=1)P(C=c))\right]}
      {-\log P(X_{\ell e}=1,C=c)},
 \quad X_{\ell e}(t)=\mathbf{1}\{e\in S_{\ell t}\}.
 \label{eq:new-npmi}
\end{equation}
Zero indicates independence; the sign indicates positive or negative association.
Degenerate joint probabilities are marked unavailable. Replacing RICE's marker
tokens with sentence-class labels is our adaptation, not its published signed
marker-ranking score. Contrasts require at least ten supporting questions in
each arm and ten with expert activation for inference. These expert-level
diagnostics and motif mining are defined by the broader interface but were not
part of the compact recovery analysis.

\paragraph{Interpretation and aggregation.}
Missing data are never encoded as zero. Question-level aggregation and
question-bootstrap intervals prevent overlapping windows from being counted
as independent observations; searched comparisons require multiple-testing
correction. Labels supplied by a retrospective judge describe observed reasoning
and are not assumed available to an online controller. These metrics supply
hypotheses about routing and reasoning dynamics, not causal claims that
persistence, instability, or any class-associated expert is helpful or harmful.

% Requires amsmath; add new_metrics to the parent document's BibTeX databases.
% Definitions describe diagnostics, not empirical results or validated controls.
\section{Temporal Routing and Reasoning-Class Dynamics}
\label{sec:new-routing-metrics}

We extend aggregate routing statistics with diagnostics of allocation persistence,
changes in mixture weights, and the organization of reasoning episodes.
The definitions below separate published foundations from adaptations introduced
for this study. They measure observed behavior; a large or small value does not,
by itself, identify useful reasoning or a beneficial intervention.

\subsection{Notation and observation scope}

At routed layer $\ell$ and reasoning-token position $t$, let
$S_{\ell t}\subseteq\{1,\ldots,E_\ell\}$ be the native selected expert set,
with $|S_{\ell t}|=k_\ell$. Let $p_{\ell t}$ denote the full router probability
vector and $a_{\ell te}\geq0$ the weight actually assigned to a selected expert.
The normalized executed mixture is
\begin{equation}
 m_{\ell te}=
 \frac{\mathbf{1}\{e\in S_{\ell t}\}\,a_{\ell te}}
      {\sum_{j\in S_{\ell t}}a_{\ell tj}}.
 \label{eq:new-executed-mixture}
\end{equation}
Thus $p_{\ell t}$ describes router preferences, whereas $m_{\ell t}$ describes
the relative weights of the selected routed experts. Neither is reconstructed
from scalar confidence summaries. Shared experts are excluded from $S_{\ell t}$
and these routed-expert distributions.

We use a trailing window $W_t^{(w)}=\{t-w+1,\ldots,t\}$ contained in one
contiguous reasoning segment. The current configuration uses
$w\in\{64,256\}$, stride $64$, and lags $h\in\{1,4,16,64\}$.
Prompt tokens, padding and gaps are excluded; insufficient history yields an
unavailable metric. All logarithms are natural, so entropies and divergences are
in nats. Layer indices are preserved throughout.

\subsection{Expert-set persistence and turnover}

Let $V_{t,h}^{(w)}$ contain positions $u\in W_t^{(w)}$ for which $u-h$ is
observed in the same reasoning segment. Lagged overlap is
\begin{equation}
 O_{\ell,h}(t;w)=\frac{1}{|V_{t,h}^{(w)}|}
 \sum_{u\in V_{t,h}^{(w)}}
 \frac{|S_{\ell u}\cap S_{\ell,u-h}|}{k_\ell},
 \qquad U_\ell(t;w)=1-O_{\ell,1}(t;w).
 \label{eq:new-overlap}
\end{equation}
Overlap lies in $[0,1]$ and scores persistence of selected identities; turnover
$U$ scores adjacent replacement. Rank changes inside an unchanged set do not
increase turnover. The temporal lag construction extends the existing adjacent
overlap metric.

\paragraph{Uniform-chance correction.}
Following the normalized overlap of Olson et al.~\cite{newmetrics_olson2025},
but applying it to temporal pairs, define
\begin{equation}
 O^{\mathrm{unif}}_{\ell,h}
 =\frac{O_{\ell,h}-k_\ell/E_\ell}{1-k_\ell/E_\ell}.
 \label{eq:new-corrected-overlap}
\end{equation}
Zero denotes the uniform independent-selection baseline; one denotes identical
sets. Negative values are possible, and the expression is undefined when
$k_\ell=E_\ell$. This affine transform is not independent evidence from raw overlap.

\paragraph{Usage-preserving order control.}
Our additional control permutes routing sets within the observed trailing
$w+h$-token pool, clipped at the segment start. If the pool contains $n$ tokens
and expert $e$ appears at $n_e$ positions, the exact expected overlap under
sampling distinct token positions is
\begin{equation}
 B_{\ell,h}=\frac{\sum_e n_e(n_e-1)}{n(n-1)k_\ell},
 \qquad O^{\mathrm{excess}}_{\ell,h}=O_{\ell,h}-B_{\ell,h}.
 \label{eq:new-shuffled-overlap}
\end{equation}
This separates temporal persistence from high overlap caused simply by repeatedly
using a small expert pool. Positive excess indicates more persistence than the
same usage frequencies predict under random ordering.

\subsection{Allocation concentration, diversity and drift}

\paragraph{Local and marginal entropy.}
LIMoE distinguishes token-local entropy from entropy of the average routing
distribution~\cite{newmetrics_mustafa2022}. We adapt this distinction from
modality-level training regularization to descriptive trailing windows.
For $r\in\{p,m\}$, set $\bar r_{\ell,W}=w^{-1}\sum_{u\in W}r_{\ell u}$ and
\begin{align}
 H^{r}_{\mathrm{local}}(W)
   &=\frac{1}{w}\sum_{u\in W}H(r_{\ell u}),
 &H^{r}_{\mathrm{marginal}}(W)&=H(\bar r_{\ell,W}),\nonumber\\
 H(v)&=-\sum_e v_e\log v_e,
 &I^{r}_{\mathrm{alloc}}(W)&=H^{r}_{\mathrm{marginal}}(W)
                           -H^{r}_{\mathrm{local}}(W).
 \label{eq:new-entropies}
\end{align}
Local entropy scores within-token dispersion. Marginal entropy scores diversity
across the window. Their nonnegative difference scores how much distributions
vary across positions: it is the mutual information between a uniformly sampled
window position and an expert drawn from that position's distribution. It is
invariant to token order, so it does not itself measure temporal drift.

For native selection counts $c_e=\sum_{u\in W}\mathbf{1}\{e\in S_{\ell u}\}$,
we additionally report $H_{\mathrm{selection}}=H((c_e/(wk_\ell))_e)$.
This describes usage diversity without using mixture weights. The corresponding
selection rates $c_e/w$ sum to $k_\ell$, not one.

\paragraph{Effective expert count.}
We express marginal executed-mixture entropy as
\begin{equation}
 N_{\mathrm{eff},\ell}(W)=\exp\!\left(H^{m}_{\mathrm{marginal}}(W)\right).
 \label{eq:new-effective-experts}
\end{equation}
It is the number of equally weighted experts with the same window-level entropy,
between $1$ and $E_\ell$. It can exceed $k_\ell$ because different tokens select
different experts; it is not the number executed per token.

\paragraph{Within-window distribution drift.}
Using Jensen--Shannon divergence~\cite{newmetrics_lin1991}, split $W$ into equal
chronological halves $W^-$ and $W^+$ and define
\begin{align}
 D^r_\ell(W)&=\operatorname{JSD}(\bar r_{\ell,W^-},\bar r_{\ell,W^+}),\nonumber\\
 \operatorname{JSD}(a,b)&=H\!\left(\frac{a+b}{2}\right)
                        -\frac{H(a)+H(b)}{2}.
 \label{eq:new-jsd}
\end{align}
The score lies in $[0,\log2]$ and measures redistribution between the two halves.
Full-router and executed-mixture drift are reported separately. Our temporal
control averages the same statistic over $B=16$ random permutations of positions
inside $W$, giving $D^{r,\mathrm{shuffle}}$ and
$D^{r,\mathrm{excess}}=D^r-D^{r,\mathrm{shuffle}}$.
Excess drift can be negative. All permutations use only the observed window.

\subsection{Native selection-boundary diagnostics}

Let $s_{\ell te}$ be the scores actually used for deterministic top-$k$
selection, including any relevant selection bias. We extend the existing
boundary-margin diagnostic with a temporal summary:
\begin{equation}
 g_{\ell t}=\min_{e\in S_{\ell t}}s_{\ell te}
            -\max_{e\notin S_{\ell t}}s_{\ell te},
 \qquad \bar g_\ell(W)=\frac{1}{w}\sum_{u\in W}g_{\ell u}.
 \label{eq:new-boundary}
\end{equation}
A small gap indicates a nearly tied selected/unselected boundary in score space;
it does not measure answer robustness. Ties give zero. This metric requires
$k_\ell<E_\ell$ and verified native score semantics; unsupported routing rules
are marked unavailable.

The study-specific weak-boundary fraction is
\begin{equation}
 F^{\mathrm{weak}}_\ell(W)=\frac{1}{w}
 \sum_{u\in W}\mathbf{1}\{g_{\ell u}\leq\theta_\ell\},
 \qquad \theta_\ell=Q_{0.10}(g_\ell;\mathcal D_{\mathrm{train}}).
 \label{eq:new-weak-boundary}
\end{equation}
The threshold is fitted only on calibration data or the training fold, counting
each token once despite overlapping windows. A larger fraction means more
near-boundary selections relative to that reference. Until such a threshold is
fitted, the fraction remains unavailable; the compact correlation report uses
the mean gap only.

\subsection{Reasoning-class dynamics}

ThinkARM motivates describing reasoning through episode sequences and transition
patterns~\cite{newmetrics_li2026}. We retain the project's seven labels
$\mathcal C=\{\mathrm{Read},\mathrm{Analyze},\mathrm{Plan},\mathrm{Implement},
\mathrm{Explore},\mathrm{Verify},\mathrm{Monitor}\}$ rather than ThinkARM's
eight-class taxonomy. Our censoring and hazard diagnostics are study-specific
extensions, not claimed reproductions of ThinkARM metrics.
Let $C_i$ be the label of reasoning sentence $i$ and $n_i$ its aligned token count.
Unknown labels break sequences rather than becoming an eighth class.

\paragraph{Occupancy, transitions and transition uncertainty.}
Sentence and token shares are respectively
$f_c=\sum_i\mathbf{1}\{C_i=c\}/N$ and
$f_c^{\mathrm{tok}}=\sum_i n_i\mathbf{1}\{C_i=c\}/\sum_i n_i$,
with denominators restricted to labeled observations and label coverage reported
separately. For adjacent labeled sentences in one contiguous segment, let $N_{cd}$
count transitions from $c$ to $d$. Then
\begin{equation}
 P_{cd}=\frac{N_{cd}}{\sum_{b\in\mathcal C}N_{cb}},\qquad
 P_{\mathrm{stay}}(c)=P_{cc},\qquad
 P_{\mathrm{exit}}(c)=1-P_{cc},\qquad
 H_{\mathrm{next}}(c)=-\sum_d P_{cd}\log P_{cd}.
 \label{eq:new-transitions}
\end{equation}
These score destination preference, persistence, switching, and uncertainty over
observed next classes. Zero source support gives an undefined rate. Natural
termination is outside this seven-class matrix; hence $1-P_{cc}$ is conditional
on an observed next labeled sentence, not an unconditional stopping probability.
Conditioning additionally on $C_{i-1}$ describes second-order transition patterns.

\paragraph{Runs, returns and motifs.}
A run is a maximal contiguous sequence of the same class. Its sentence/token
length scores sustained allocation to that reasoning activity; its current age
uses only the observed prefix. Entry after a missing-label gap is unobserved, so
full duration is unavailable. A gap, cap or unknown ending right-censors the run;
observed natural termination is a competing exit event.
Return intervals measure the sentence/token distance between successive run
starts of the same class, without crossing gaps. Counts of contiguous label
bigrams and trigrams describe recurring local motifs. Motif counts alone are
not significance scores and do not establish purposeful self-correction.

\paragraph{Destination-specific exit hazards.}
For runs with observed entry, let $Y_c(a)$ count runs whose outgoing boundary at
sentence age $a$ is observable, and $D_{cd}(a)$ count exits to destination $d$ at
that age. Our discrete hazard is
\begin{equation}
 h_{c\rightarrow d}(a)=D_{cd}(a)/Y_c(a),
 \qquad d\in(\mathcal C\setminus\{c\})\cup\{\mathrm{natural\ termination}\}.
 \label{eq:new-hazard}
\end{equation}
It scores the chance of a particular observed exit at that age, conditional on
reaching an observable boundary. Right-censored runs contribute earlier
boundaries but not their unobserved final boundary. Empty risk sets remain
undefined; observed run lengths must not be mistaken for uncensored dwell times.

\subsection{Relating routing to reasoning classes}

\paragraph{Within-question class contrasts.}
For a scalar window metric $z$, let $\bar z_{q,c}$ average windows ending in
class $c$ within question $q$, and $\bar z_{q,\neg c}$ average other labeled
endpoints. Our contrast is
\begin{equation}
 \Delta_c(z)=\frac{1}{|\mathcal Q_c|}
 \sum_{q\in\mathcal Q_c}(\bar z_{q,c}-\bar z_{q,\neg c}),
 \label{eq:new-class-contrast}
\end{equation}
where $\mathcal Q_c$ contains questions with both kinds of window. Layers are
averaged equally within each question. Positive values indicate larger metric
values around that class. A trailing window can include earlier classes, so this
is an endpoint-conditioned association. The saved original class-view comparison
instead uses class minus all reasoning and is kept as a distinct estimand.

\paragraph{Expert-specific association diagnostics.}
The broader diagnostic interface also compares expert activation rates, adapting
SteerMoE's contrast principle~\cite{newmetrics_fayyaz2026} from paired behavioral
prompts to observed classes. For tokens with label $c$, define
\begin{equation}
 R_{\ell e\mid c}=\frac{\sum_t\mathbf{1}\{C(t)=c,\ e\in S_{\ell t}\}}
                          {\sum_t\mathbf{1}\{C(t)=c\}},
 \qquad A_{\ell e,c}=R_{\ell e\mid c}-R_{\ell e\mid\neg c}.
 \label{eq:new-activation-contrast}
\end{equation}
This is a token-weighted selection-rate difference; positive values indicate
class-associated recruitment. Expert identity always includes its layer.
A secondary normalized pointwise mutual information score, following the
association measure used by RICE~\cite{newmetrics_wang2025}, is
\begin{equation}
 \operatorname{nPMI}(X_{\ell e},c)=
 \frac{\log\!\left[P(X_{\ell e}=1,C=c)/
                (P(X_{\ell e}=1)P(C=c))\right]}
      {-\log P(X_{\ell e}=1,C=c)},
 \quad X_{\ell e}(t)=\mathbf{1}\{e\in S_{\ell t}\}.
 \label{eq:new-npmi}
\end{equation}
Zero indicates independence; the sign indicates positive or negative association.
Degenerate joint probabilities are marked unavailable. Replacing RICE's marker
tokens with sentence-class labels is our adaptation, not its published signed
marker-ranking score. Contrasts require at least ten supporting questions in
each arm and ten with expert activation for inference. These expert-level
diagnostics and motif mining are defined by the broader interface but were not
part of the compact recovery analysis.

\paragraph{Interpretation and aggregation.}
Missing data are never encoded as zero. Question-level aggregation and
question-bootstrap intervals prevent overlapping windows from being counted
as independent observations; searched comparisons require multiple-testing
correction. Labels supplied by a retrospective judge describe observed reasoning
and are not assumed available to an online controller. These metrics supply
hypotheses about routing and reasoning dynamics, not causal claims that
persistence, instability, or any class-associated expert is helpful or harmful.

% Requires amsmath; add new_metrics to the parent document's BibTeX databases.
% Definitions describe diagnostics, not empirical results or validated controls.
\section{Temporal Routing and Reasoning-Class Dynamics}
\label{sec:new-routing-metrics}

We extend aggregate routing statistics with diagnostics of allocation persistence,
changes in mixture weights, and the organization of reasoning episodes.
The definitions below separate published foundations from adaptations introduced
for this study. They measure observed behavior; a large or small value does not,
by itself, identify useful reasoning or a beneficial intervention.

\subsection{Notation and observation scope}

At routed layer $\ell$ and reasoning-token position $t$, let
$S_{\ell t}\subseteq\{1,\ldots,E_\ell\}$ be the native selected expert set,
with $|S_{\ell t}|=k_\ell$. Let $p_{\ell t}$ denote the full router probability
vector and $a_{\ell te}\geq0$ the weight actually assigned to a selected expert.
The normalized executed mixture is
\begin{equation}
 m_{\ell te}=
 \frac{\mathbf{1}\{e\in S_{\ell t}\}\,a_{\ell te}}
      {\sum_{j\in S_{\ell t}}a_{\ell tj}}.
 \label{eq:new-executed-mixture}
\end{equation}
Thus $p_{\ell t}$ describes router preferences, whereas $m_{\ell t}$ describes
the relative weights of the selected routed experts. Neither is reconstructed
from scalar confidence summaries. Shared experts are excluded from $S_{\ell t}$
and these routed-expert distributions.

We use a trailing window $W_t^{(w)}=\{t-w+1,\ldots,t\}$ contained in one
contiguous reasoning segment. The current configuration uses
$w\in\{64,256\}$, stride $64$, and lags $h\in\{1,4,16,64\}$.
Prompt tokens, padding and gaps are excluded; insufficient history yields an
unavailable metric. All logarithms are natural, so entropies and divergences are
in nats. Layer indices are preserved throughout.

\subsection{Expert-set persistence and turnover}

Let $V_{t,h}^{(w)}$ contain positions $u\in W_t^{(w)}$ for which $u-h$ is
observed in the same reasoning segment. Lagged overlap is
\begin{equation}
 O_{\ell,h}(t;w)=\frac{1}{|V_{t,h}^{(w)}|}
 \sum_{u\in V_{t,h}^{(w)}}
 \frac{|S_{\ell u}\cap S_{\ell,u-h}|}{k_\ell},
 \qquad U_\ell(t;w)=1-O_{\ell,1}(t;w).
 \label{eq:new-overlap}
\end{equation}
Overlap lies in $[0,1]$ and scores persistence of selected identities; turnover
$U$ scores adjacent replacement. Rank changes inside an unchanged set do not
increase turnover. The temporal lag construction extends the existing adjacent
overlap metric.

\paragraph{Uniform-chance correction.}
Following the normalized overlap of Olson et al.~\cite{newmetrics_olson2025},
but applying it to temporal pairs, define
\begin{equation}
 O^{\mathrm{unif}}_{\ell,h}
 =\frac{O_{\ell,h}-k_\ell/E_\ell}{1-k_\ell/E_\ell}.
 \label{eq:new-corrected-overlap}
\end{equation}
Zero denotes the uniform independent-selection baseline; one denotes identical
sets. Negative values are possible, and the expression is undefined when
$k_\ell=E_\ell$. This affine transform is not independent evidence from raw overlap.

\paragraph{Usage-preserving order control.}
Our additional control permutes routing sets within the observed trailing
$w+h$-token pool, clipped at the segment start. If the pool contains $n$ tokens
and expert $e$ appears at $n_e$ positions, the exact expected overlap under
sampling distinct token positions is
\begin{equation}
 B_{\ell,h}=\frac{\sum_e n_e(n_e-1)}{n(n-1)k_\ell},
 \qquad O^{\mathrm{excess}}_{\ell,h}=O_{\ell,h}-B_{\ell,h}.
 \label{eq:new-shuffled-overlap}
\end{equation}
This separates temporal persistence from high overlap caused simply by repeatedly
using a small expert pool. Positive excess indicates more persistence than the
same usage frequencies predict under random ordering.

\subsection{Allocation concentration, diversity and drift}

\paragraph{Local and marginal entropy.}
LIMoE distinguishes token-local entropy from entropy of the average routing
distribution~\cite{newmetrics_mustafa2022}. We adapt this distinction from
modality-level training regularization to descriptive trailing windows.
For $r\in\{p,m\}$, set $\bar r_{\ell,W}=w^{-1}\sum_{u\in W}r_{\ell u}$ and
\begin{align}
 H^{r}_{\mathrm{local}}(W)
   &=\frac{1}{w}\sum_{u\in W}H(r_{\ell u}),
 &H^{r}_{\mathrm{marginal}}(W)&=H(\bar r_{\ell,W}),\nonumber\\
 H(v)&=-\sum_e v_e\log v_e,
 &I^{r}_{\mathrm{alloc}}(W)&=H^{r}_{\mathrm{marginal}}(W)
                           -H^{r}_{\mathrm{local}}(W).
 \label{eq:new-entropies}
\end{align}
Local entropy scores within-token dispersion. Marginal entropy scores diversity
across the window. Their nonnegative difference scores how much distributions
vary across positions: it is the mutual information between a uniformly sampled
window position and an expert drawn from that position's distribution. It is
invariant to token order, so it does not itself measure temporal drift.

For native selection counts $c_e=\sum_{u\in W}\mathbf{1}\{e\in S_{\ell u}\}$,
we additionally report $H_{\mathrm{selection}}=H((c_e/(wk_\ell))_e)$.
This describes usage diversity without using mixture weights. The corresponding
selection rates $c_e/w$ sum to $k_\ell$, not one.

\paragraph{Effective expert count.}
We express marginal executed-mixture entropy as
\begin{equation}
 N_{\mathrm{eff},\ell}(W)=\exp\!\left(H^{m}_{\mathrm{marginal}}(W)\right).
 \label{eq:new-effective-experts}
\end{equation}
It is the number of equally weighted experts with the same window-level entropy,
between $1$ and $E_\ell$. It can exceed $k_\ell$ because different tokens select
different experts; it is not the number executed per token.

\paragraph{Within-window distribution drift.}
Using Jensen--Shannon divergence~\cite{newmetrics_lin1991}, split $W$ into equal
chronological halves $W^-$ and $W^+$ and define
\begin{align}
 D^r_\ell(W)&=\operatorname{JSD}(\bar r_{\ell,W^-},\bar r_{\ell,W^+}),\nonumber\\
 \operatorname{JSD}(a,b)&=H\!\left(\frac{a+b}{2}\right)
                        -\frac{H(a)+H(b)}{2}.
 \label{eq:new-jsd}
\end{align}
The score lies in $[0,\log2]$ and measures redistribution between the two halves.
Full-router and executed-mixture drift are reported separately. Our temporal
control averages the same statistic over $B=16$ random permutations of positions
inside $W$, giving $D^{r,\mathrm{shuffle}}$ and
$D^{r,\mathrm{excess}}=D^r-D^{r,\mathrm{shuffle}}$.
Excess drift can be negative. All permutations use only the observed window.

\subsection{Native selection-boundary diagnostics}

Let $s_{\ell te}$ be the scores actually used for deterministic top-$k$
selection, including any relevant selection bias. We extend the existing
boundary-margin diagnostic with a temporal summary:
\begin{equation}
 g_{\ell t}=\min_{e\in S_{\ell t}}s_{\ell te}
            -\max_{e\notin S_{\ell t}}s_{\ell te},
 \qquad \bar g_\ell(W)=\frac{1}{w}\sum_{u\in W}g_{\ell u}.
 \label{eq:new-boundary}
\end{equation}
A small gap indicates a nearly tied selected/unselected boundary in score space;
it does not measure answer robustness. Ties give zero. This metric requires
$k_\ell<E_\ell$ and verified native score semantics; unsupported routing rules
are marked unavailable.

The study-specific weak-boundary fraction is
\begin{equation}
 F^{\mathrm{weak}}_\ell(W)=\frac{1}{w}
 \sum_{u\in W}\mathbf{1}\{g_{\ell u}\leq\theta_\ell\},
 \qquad \theta_\ell=Q_{0.10}(g_\ell;\mathcal D_{\mathrm{train}}).
 \label{eq:new-weak-boundary}
\end{equation}
The threshold is fitted only on calibration data or the training fold, counting
each token once despite overlapping windows. A larger fraction means more
near-boundary selections relative to that reference. Until such a threshold is
fitted, the fraction remains unavailable; the compact correlation report uses
the mean gap only.

\subsection{Reasoning-class dynamics}

ThinkARM motivates describing reasoning through episode sequences and transition
patterns~\cite{newmetrics_li2026}. We retain the project's seven labels
$\mathcal C=\{\mathrm{Read},\mathrm{Analyze},\mathrm{Plan},\mathrm{Implement},
\mathrm{Explore},\mathrm{Verify},\mathrm{Monitor}\}$ rather than ThinkARM's
eight-class taxonomy. Our censoring and hazard diagnostics are study-specific
extensions, not claimed reproductions of ThinkARM metrics.
Let $C_i$ be the label of reasoning sentence $i$ and $n_i$ its aligned token count.
Unknown labels break sequences rather than becoming an eighth class.

\paragraph{Occupancy, transitions and transition uncertainty.}
Sentence and token shares are respectively
$f_c=\sum_i\mathbf{1}\{C_i=c\}/N$ and
$f_c^{\mathrm{tok}}=\sum_i n_i\mathbf{1}\{C_i=c\}/\sum_i n_i$,
with denominators restricted to labeled observations and label coverage reported
separately. For adjacent labeled sentences in one contiguous segment, let $N_{cd}$
count transitions from $c$ to $d$. Then
\begin{equation}
 P_{cd}=\frac{N_{cd}}{\sum_{b\in\mathcal C}N_{cb}},\qquad
 P_{\mathrm{stay}}(c)=P_{cc},\qquad
 P_{\mathrm{exit}}(c)=1-P_{cc},\qquad
 H_{\mathrm{next}}(c)=-\sum_d P_{cd}\log P_{cd}.
 \label{eq:new-transitions}
\end{equation}
These score destination preference, persistence, switching, and uncertainty over
observed next classes. Zero source support gives an undefined rate. Natural
termination is outside this seven-class matrix; hence $1-P_{cc}$ is conditional
on an observed next labeled sentence, not an unconditional stopping probability.
Conditioning additionally on $C_{i-1}$ describes second-order transition patterns.

\paragraph{Runs, returns and motifs.}
A run is a maximal contiguous sequence of the same class. Its sentence/token
length scores sustained allocation to that reasoning activity; its current age
uses only the observed prefix. Entry after a missing-label gap is unobserved, so
full duration is unavailable. A gap, cap or unknown ending right-censors the run;
observed natural termination is a competing exit event.
Return intervals measure the sentence/token distance between successive run
starts of the same class, without crossing gaps. Counts of contiguous label
bigrams and trigrams describe recurring local motifs. Motif counts alone are
not significance scores and do not establish purposeful self-correction.

\paragraph{Destination-specific exit hazards.}
For runs with observed entry, let $Y_c(a)$ count runs whose outgoing boundary at
sentence age $a$ is observable, and $D_{cd}(a)$ count exits to destination $d$ at
that age. Our discrete hazard is
\begin{equation}
 h_{c\rightarrow d}(a)=D_{cd}(a)/Y_c(a),
 \qquad d\in(\mathcal C\setminus\{c\})\cup\{\mathrm{natural\ termination}\}.
 \label{eq:new-hazard}
\end{equation}
It scores the chance of a particular observed exit at that age, conditional on
reaching an observable boundary. Right-censored runs contribute earlier
boundaries but not their unobserved final boundary. Empty risk sets remain
undefined; observed run lengths must not be mistaken for uncensored dwell times.

\subsection{Relating routing to reasoning classes}

\paragraph{Within-question class contrasts.}
For a scalar window metric $z$, let $\bar z_{q,c}$ average windows ending in
class $c$ within question $q$, and $\bar z_{q,\neg c}$ average other labeled
endpoints. Our contrast is
\begin{equation}
 \Delta_c(z)=\frac{1}{|\mathcal Q_c|}
 \sum_{q\in\mathcal Q_c}(\bar z_{q,c}-\bar z_{q,\neg c}),
 \label{eq:new-class-contrast}
\end{equation}
where $\mathcal Q_c$ contains questions with both kinds of window. Layers are
averaged equally within each question. Positive values indicate larger metric
values around that class. A trailing window can include earlier classes, so this
is an endpoint-conditioned association. The saved original class-view comparison
instead uses class minus all reasoning and is kept as a distinct estimand.

\paragraph{Expert-specific association diagnostics.}
The broader diagnostic interface also compares expert activation rates, adapting
SteerMoE's contrast principle~\cite{newmetrics_fayyaz2026} from paired behavioral
prompts to observed classes. For tokens with label $c$, define
\begin{equation}
 R_{\ell e\mid c}=\frac{\sum_t\mathbf{1}\{C(t)=c,\ e\in S_{\ell t}\}}
                          {\sum_t\mathbf{1}\{C(t)=c\}},
 \qquad A_{\ell e,c}=R_{\ell e\mid c}-R_{\ell e\mid\neg c}.
 \label{eq:new-activation-contrast}
\end{equation}
This is a token-weighted selection-rate difference; positive values indicate
class-associated recruitment. Expert identity always includes its layer.
A secondary normalized pointwise mutual information score, following the
association measure used by RICE~\cite{newmetrics_wang2025}, is
\begin{equation}
 \operatorname{nPMI}(X_{\ell e},c)=
 \frac{\log\!\left[P(X_{\ell e}=1,C=c)/
                (P(X_{\ell e}=1)P(C=c))\right]}
      {-\log P(X_{\ell e}=1,C=c)},
 \quad X_{\ell e}(t)=\mathbf{1}\{e\in S_{\ell t}\}.
 \label{eq:new-npmi}
\end{equation}
Zero indicates independence; the sign indicates positive or negative association.
Degenerate joint probabilities are marked unavailable. Replacing RICE's marker
tokens with sentence-class labels is our adaptation, not its published signed
marker-ranking score. Contrasts require at least ten supporting questions in
each arm and ten with expert activation for inference. These expert-level
diagnostics and motif mining are defined by the broader interface but were not
part of the compact recovery analysis.

\paragraph{Interpretation and aggregation.}
Missing data are never encoded as zero. Question-level aggregation and
question-bootstrap intervals prevent overlapping windows from being counted
as independent observations; searched comparisons require multiple-testing
correction. Labels supplied by a retrospective judge describe observed reasoning
and are not assumed available to an online controller. These metrics supply
hypotheses about routing and reasoning dynamics, not causal claims that
persistence, instability, or any class-associated expert is helpful or harmful.

% Requires amsmath; add new_metrics to the parent document's BibTeX databases.
% Definitions describe diagnostics, not empirical results or validated controls.
\section{Temporal Routing and Reasoning-Class Dynamics}
\label{sec:new-routing-metrics}

We extend aggregate routing statistics with diagnostics of allocation persistence,
changes in mixture weights, and the organization of reasoning episodes.
The definitions below separate published foundations from adaptations introduced
for this study. They measure observed behavior; a large or small value does not,
by itself, identify useful reasoning or a beneficial intervention.

\subsection{Notation and observation scope}

At routed layer $\ell$ and reasoning-token position $t$, let
$S_{\ell t}\subseteq\{1,\ldots,E_\ell\}$ be the native selected expert set,
with $|S_{\ell t}|=k_\ell$. Let $p_{\ell t}$ denote the full router probability
vector and $a_{\ell te}\geq0$ the weight actually assigned to a selected expert.
The normalized executed mixture is
\begin{equation}
 m_{\ell te}=
 \frac{\mathbf{1}\{e\in S_{\ell t}\}\,a_{\ell te}}
      {\sum_{j\in S_{\ell t}}a_{\ell tj}}.
 \label{eq:new-executed-mixture}
\end{equation}
Thus $p_{\ell t}$ describes router preferences, whereas $m_{\ell t}$ describes
the relative weights of the selected routed experts. Neither is reconstructed
from scalar confidence summaries. Shared experts are excluded from $S_{\ell t}$
and these routed-expert distributions.

We use a trailing window $W_t^{(w)}=\{t-w+1,\ldots,t\}$ contained in one
contiguous reasoning segment. The current configuration uses
$w\in\{64,256\}$, stride $64$, and lags $h\in\{1,4,16,64\}$.
Prompt tokens, padding and gaps are excluded; insufficient history yields an
unavailable metric. All logarithms are natural, so entropies and divergences are
in nats. Layer indices are preserved throughout.

\subsection{Expert-set persistence and turnover}

Let $V_{t,h}^{(w)}$ contain positions $u\in W_t^{(w)}$ for which $u-h$ is
observed in the same reasoning segment. Lagged overlap is
\begin{equation}
 O_{\ell,h}(t;w)=\frac{1}{|V_{t,h}^{(w)}|}
 \sum_{u\in V_{t,h}^{(w)}}
 \frac{|S_{\ell u}\cap S_{\ell,u-h}|}{k_\ell},
 \qquad U_\ell(t;w)=1-O_{\ell,1}(t;w).
 \label{eq:new-overlap}
\end{equation}
Overlap lies in $[0,1]$ and scores persistence of selected identities; turnover
$U$ scores adjacent replacement. Rank changes inside an unchanged set do not
increase turnover. The temporal lag construction extends the existing adjacent
overlap metric.

\paragraph{Uniform-chance correction.}
Following the normalized overlap of Olson et al.~\cite{newmetrics_olson2025},
but applying it to temporal pairs, define
\begin{equation}
 O^{\mathrm{unif}}_{\ell,h}
 =\frac{O_{\ell,h}-k_\ell/E_\ell}{1-k_\ell/E_\ell}.
 \label{eq:new-corrected-overlap}
\end{equation}
Zero denotes the uniform independent-selection baseline; one denotes identical
sets. Negative values are possible, and the expression is undefined when
$k_\ell=E_\ell$. This affine transform is not independent evidence from raw overlap.

\paragraph{Usage-preserving order control.}
Our additional control permutes routing sets within the observed trailing
$w+h$-token pool, clipped at the segment start. If the pool contains $n$ tokens
and expert $e$ appears at $n_e$ positions, the exact expected overlap under
sampling distinct token positions is
\begin{equation}
 B_{\ell,h}=\frac{\sum_e n_e(n_e-1)}{n(n-1)k_\ell},
 \qquad O^{\mathrm{excess}}_{\ell,h}=O_{\ell,h}-B_{\ell,h}.
 \label{eq:new-shuffled-overlap}
\end{equation}
This separates temporal persistence from high overlap caused simply by repeatedly
using a small expert pool. Positive excess indicates more persistence than the
same usage frequencies predict under random ordering.

\subsection{Allocation concentration, diversity and drift}

\paragraph{Local and marginal entropy.}
LIMoE distinguishes token-local entropy from entropy of the average routing
distribution~\cite{newmetrics_mustafa2022}. We adapt this distinction from
modality-level training regularization to descriptive trailing windows.
For $r\in\{p,m\}$, set $\bar r_{\ell,W}=w^{-1}\sum_{u\in W}r_{\ell u}$ and
\begin{align}
 H^{r}_{\mathrm{local}}(W)
   &=\frac{1}{w}\sum_{u\in W}H(r_{\ell u}),
 &H^{r}_{\mathrm{marginal}}(W)&=H(\bar r_{\ell,W}),\nonumber\\
 H(v)&=-\sum_e v_e\log v_e,
 &I^{r}_{\mathrm{alloc}}(W)&=H^{r}_{\mathrm{marginal}}(W)
                           -H^{r}_{\mathrm{local}}(W).
 \label{eq:new-entropies}
\end{align}
Local entropy scores within-token dispersion. Marginal entropy scores diversity
across the window. Their nonnegative difference scores how much distributions
vary across positions: it is the mutual information between a uniformly sampled
window position and an expert drawn from that position's distribution. It is
invariant to token order, so it does not itself measure temporal drift.

For native selection counts $c_e=\sum_{u\in W}\mathbf{1}\{e\in S_{\ell u}\}$,
we additionally report $H_{\mathrm{selection}}=H((c_e/(wk_\ell))_e)$.
This describes usage diversity without using mixture weights. The corresponding
selection rates $c_e/w$ sum to $k_\ell$, not one.

\paragraph{Effective expert count.}
We express marginal executed-mixture entropy as
\begin{equation}
 N_{\mathrm{eff},\ell}(W)=\exp\!\left(H^{m}_{\mathrm{marginal}}(W)\right).
 \label{eq:new-effective-experts}
\end{equation}
It is the number of equally weighted experts with the same window-level entropy,
between $1$ and $E_\ell$. It can exceed $k_\ell$ because different tokens select
different experts; it is not the number executed per token.

\paragraph{Within-window distribution drift.}
Using Jensen--Shannon divergence~\cite{newmetrics_lin1991}, split $W$ into equal
chronological halves $W^-$ and $W^+$ and define
\begin{align}
 D^r_\ell(W)&=\operatorname{JSD}(\bar r_{\ell,W^-},\bar r_{\ell,W^+}),\nonumber\\
 \operatorname{JSD}(a,b)&=H\!\left(\frac{a+b}{2}\right)
                        -\frac{H(a)+H(b)}{2}.
 \label{eq:new-jsd}
\end{align}
The score lies in $[0,\log2]$ and measures redistribution between the two halves.
Full-router and executed-mixture drift are reported separately. Our temporal
control averages the same statistic over $B=16$ random permutations of positions
inside $W$, giving $D^{r,\mathrm{shuffle}}$ and
$D^{r,\mathrm{excess}}=D^r-D^{r,\mathrm{shuffle}}$.
Excess drift can be negative. All permutations use only the observed window.

\subsection{Native selection-boundary diagnostics}

Let $s_{\ell te}$ be the scores actually used for deterministic top-$k$
selection, including any relevant selection bias. We extend the existing
boundary-margin diagnostic with a temporal summary:
\begin{equation}
 g_{\ell t}=\min_{e\in S_{\ell t}}s_{\ell te}
            -\max_{e\notin S_{\ell t}}s_{\ell te},
 \qquad \bar g_\ell(W)=\frac{1}{w}\sum_{u\in W}g_{\ell u}.
 \label{eq:new-boundary}
\end{equation}
A small gap indicates a nearly tied selected/unselected boundary in score space;
it does not measure answer robustness. Ties give zero. This metric requires
$k_\ell<E_\ell$ and verified native score semantics; unsupported routing rules
are marked unavailable.

The study-specific weak-boundary fraction is
\begin{equation}
 F^{\mathrm{weak}}_\ell(W)=\frac{1}{w}
 \sum_{u\in W}\mathbf{1}\{g_{\ell u}\leq\theta_\ell\},
 \qquad \theta_\ell=Q_{0.10}(g_\ell;\mathcal D_{\mathrm{train}}).
 \label{eq:new-weak-boundary}
\end{equation}
The threshold is fitted only on calibration data or the training fold, counting
each token once despite overlapping windows. A larger fraction means more
near-boundary selections relative to that reference. Until such a threshold is
fitted, the fraction remains unavailable; the compact correlation report uses
the mean gap only.

\subsection{Reasoning-class dynamics}

ThinkARM motivates describing reasoning through episode sequences and transition
patterns~\cite{newmetrics_li2026}. We retain the project's seven labels
$\mathcal C=\{\mathrm{Read},\mathrm{Analyze},\mathrm{Plan},\mathrm{Implement},
\mathrm{Explore},\mathrm{Verify},\mathrm{Monitor}\}$ rather than ThinkARM's
eight-class taxonomy. Our censoring and hazard diagnostics are study-specific
extensions, not claimed reproductions of ThinkARM metrics.
Let $C_i$ be the label of reasoning sentence $i$ and $n_i$ its aligned token count.
Unknown labels break sequences rather than becoming an eighth class.

\paragraph{Occupancy, transitions and transition uncertainty.}
Sentence and token shares are respectively
$f_c=\sum_i\mathbf{1}\{C_i=c\}/N$ and
$f_c^{\mathrm{tok}}=\sum_i n_i\mathbf{1}\{C_i=c\}/\sum_i n_i$,
with denominators restricted to labeled observations and label coverage reported
separately. For adjacent labeled sentences in one contiguous segment, let $N_{cd}$
count transitions from $c$ to $d$. Then
\begin{equation}
 P_{cd}=\frac{N_{cd}}{\sum_{b\in\mathcal C}N_{cb}},\qquad
 P_{\mathrm{stay}}(c)=P_{cc},\qquad
 P_{\mathrm{exit}}(c)=1-P_{cc},\qquad
 H_{\mathrm{next}}(c)=-\sum_d P_{cd}\log P_{cd}.
 \label{eq:new-transitions}
\end{equation}
These score destination preference, persistence, switching, and uncertainty over
observed next classes. Zero source support gives an undefined rate. Natural
termination is outside this seven-class matrix; hence $1-P_{cc}$ is conditional
on an observed next labeled sentence, not an unconditional stopping probability.
Conditioning additionally on $C_{i-1}$ describes second-order transition patterns.

\paragraph{Runs, returns and motifs.}
A run is a maximal contiguous sequence of the same class. Its sentence/token
length scores sustained allocation to that reasoning activity; its current age
uses only the observed prefix. Entry after a missing-label gap is unobserved, so
full duration is unavailable. A gap, cap or unknown ending right-censors the run;
observed natural termination is a competing exit event.
Return intervals measure the sentence/token distance between successive run
starts of the same class, without crossing gaps. Counts of contiguous label
bigrams and trigrams describe recurring local motifs. Motif counts alone are
not significance scores and do not establish purposeful self-correction.

\paragraph{Destination-specific exit hazards.}
For runs with observed entry, let $Y_c(a)$ count runs whose outgoing boundary at
sentence age $a$ is observable, and $D_{cd}(a)$ count exits to destination $d$ at
that age. Our discrete hazard is
\begin{equation}
 h_{c\rightarrow d}(a)=D_{cd}(a)/Y_c(a),
 \qquad d\in(\mathcal C\setminus\{c\})\cup\{\mathrm{natural\ termination}\}.
 \label{eq:new-hazard}
\end{equation}
It scores the chance of a particular observed exit at that age, conditional on
reaching an observable boundary. Right-censored runs contribute earlier
boundaries but not their unobserved final boundary. Empty risk sets remain
undefined; observed run lengths must not be mistaken for uncensored dwell times.

\subsection{Relating routing to reasoning classes}

\paragraph{Within-question class contrasts.}
For a scalar window metric $z$, let $\bar z_{q,c}$ average windows ending in
class $c$ within question $q$, and $\bar z_{q,\neg c}$ average other labeled
endpoints. Our contrast is
\begin{equation}
 \Delta_c(z)=\frac{1}{|\mathcal Q_c|}
 \sum_{q\in\mathcal Q_c}(\bar z_{q,c}-\bar z_{q,\neg c}),
 \label{eq:new-class-contrast}
\end{equation}
where $\mathcal Q_c$ contains questions with both kinds of window. Layers are
averaged equally within each question. Positive values indicate larger metric
values around that class. A trailing window can include earlier classes, so this
is an endpoint-conditioned association. The saved original class-view comparison
instead uses class minus all reasoning and is kept as a distinct estimand.

\paragraph{Expert-specific association diagnostics.}
The broader diagnostic interface also compares expert activation rates, adapting
SteerMoE's contrast principle~\cite{newmetrics_fayyaz2026} from paired behavioral
prompts to observed classes. For tokens with label $c$, define
\begin{equation}
 R_{\ell e\mid c}=\frac{\sum_t\mathbf{1}\{C(t)=c,\ e\in S_{\ell t}\}}
                          {\sum_t\mathbf{1}\{C(t)=c\}},
 \qquad A_{\ell e,c}=R_{\ell e\mid c}-R_{\ell e\mid\neg c}.
 \label{eq:new-activation-contrast}
\end{equation}
This is a token-weighted selection-rate difference; positive values indicate
class-associated recruitment. Expert identity always includes its layer.
A secondary normalized pointwise mutual information score, following the
association measure used by RICE~\cite{newmetrics_wang2025}, is
\begin{equation}
 \operatorname{nPMI}(X_{\ell e},c)=
 \frac{\log\!\left[P(X_{\ell e}=1,C=c)/
                (P(X_{\ell e}=1)P(C=c))\right]}
      {-\log P(X_{\ell e}=1,C=c)},
 \quad X_{\ell e}(t)=\mathbf{1}\{e\in S_{\ell t}\}.
 \label{eq:new-npmi}
\end{equation}
Zero indicates independence; the sign indicates positive or negative association.
Degenerate joint probabilities are marked unavailable. Replacing RICE's marker
tokens with sentence-class labels is our adaptation, not its published signed
marker-ranking score. Contrasts require at least ten supporting questions in
each arm and ten with expert activation for inference. These expert-level
diagnostics and motif mining are defined by the broader interface but were not
part of the compact recovery analysis.

\paragraph{Interpretation and aggregation.}
Missing data are never encoded as zero. Question-level aggregation and
question-bootstrap intervals prevent overlapping windows from being counted
as independent observations; searched comparisons require multiple-testing
correction. Labels supplied by a retrospective judge describe observed reasoning
and are not assumed available to an online controller. These metrics supply
hypotheses about routing and reasoning dynamics, not causal claims that
persistence, instability, or any class-associated expert is helpful or harmful.
